\chapter{Hashtables: Open addressing}\label{chap:hashtables:open}

Open addressing handles collisions by probing for the next available slot within the array itself. All elements are stored directly in the bucket array.

\section{Design rationale}

We implemented open addressing with linear probing. When a collision occurs, we check the next slot, wrapping around if necessary. This approach offers better cache locality than closed addressing since elements are stored contiguously.

\section{Example variant: Tombstone reuse}

When an element is removed, we cannot simply mark the slot as empty, as this would break the probe chain for other elements. A common solution is to use a "tombstone" marker.

In our implementation, we support tombstone reuse. When inserting a new element, if we encounter a tombstone during probing, we can overwrite it with the new element (provided the key doesn't already exist further down the chain). This helps keep the table from filling up with tombstones and reduces the need for frequent rehashing.

\section{Example variant: Backward-shift deletion (no tombstones)}

To avoid the accumulation of tombstones entirely, we also implemented backward-shift deletion. When an element is removed, we check subsequent elements in the probe chain. If an element is found that "belongs" to a slot prior to the empty slot (due to wrapping or previous collisions), we shift it back into the empty slot. This process repeats until we find an empty slot or an element that is in its correct home position.

This approach keeps the table clean and ensures that search times remain optimal, at the cost of slightly more complex deletion logic.

\section{Example variant: Robin-Hood hashing + Backward-shift deletion}

We also explored Robin-Hood hashing, a technique that minimizes the variance of probe lengths. During insertion, if we encounter an element that is "richer" (closer to its home slot) than the element we are inserting, we swap them. The inserted element "steals" the slot, and the displaced element continues to probe for a new spot. This ensures that no element is too far from its home slot, leading to more predictable performance. Combined with backward-shift deletion, this results in a highly efficient and robust hashtable.
